\subsection{ন্যানো টেকনোলজি}

\begin{learningbox}
এই অংশ শেষে শিক্ষার্থী পারবে:
\begin{itemize}
\item ন্যানো টেকনোলজির ধারণা ব্যাখ্যা করতে।
\item nanometer, nanomaterial ও nanoscale-এর সম্পর্ক বুঝতে।
\item চিকিৎসা, ইলেকট্রনিক্স, কৃষি ও শিল্পে ন্যানো টেকনোলজির ব্যবহার লিখতে।
\item ন্যানো টেকনোলজির সুবিধা ও সম্ভাব্য ঝুঁকি বিশ্লেষণ করতে।
\end{itemize}
\end{learningbox}

\subsubsection{ন্যানো টেকনোলজির ধারণা}

Nano অর্থ অত্যন্ত ক্ষুদ্র। ন্যানো টেকনোলজি হলো nanoscale-এ পদার্থের গঠন, বৈশিষ্ট্য ও ব্যবহার নিয়ে কাজ করার প্রযুক্তি। National Nanotechnology Initiative অনুযায়ী nanotechnology সাধারণত 1 থেকে 100 nanometer মাত্রায় পদার্থ নিয়ন্ত্রণ ও ব্যবহার নিয়ে কাজ করে।

\begin{definitionbox}{সংজ্ঞা}
ন্যানো টেকনোলজি হলো 1-100 nanometer মাত্রায় পদার্থের গঠন ও বৈশিষ্ট্য নিয়ন্ত্রণ করে নতুন material, device বা system তৈরি করার প্রযুক্তি।
\end{definitionbox}

একই পদার্থ nanoscale-এ গেলে তার রং, শক্তি, chemical reactivity, electrical property বা surface area ভিন্ন হতে পারে। এই বিশেষ বৈশিষ্ট্যই nanotechnology-কে গুরুত্বপূর্ণ করে।

\subsubsection{ন্যানোমিটার ও ন্যানো উপাদান}

১ nanometer হলো ১ meter-এর এক বিলিয়ন ভাগের এক ভাগ। মানুষের চুলের ব্যাস প্রায় হাজার হাজার nanometer; তাই nanoscale অত্যন্ত ক্ষুদ্র।

\begin{keypointbox}{ন্যানো উপাদানের উদাহরণ}
\begin{itemize}
\item nanoparticle।
\item carbon nanotube।
\item graphene।
\item quantum dot।
\item nano coating।
\item nano sensor।
\end{itemize}
\end{keypointbox}

\subsubsection{চিকিৎসায় ব্যবহার}

Nanotechnology medicine-এ drug delivery, cancer therapy research, biosensor, diagnostic imaging এবং tissue engineering-এ ব্যবহার করা হয়। Nanoparticle দিয়ে নির্দিষ্ট cell বা tissue-তে drug পৌঁছানোর গবেষণা চলছে।

\begin{examplebox}{Targeted drug delivery}
সাধারণ ওষুধ পুরো শরীরে ছড়িয়ে পড়ে side effect তৈরি করতে পারে। nanoscale carrier ব্যবহার করে নির্দিষ্ট tissue-তে ওষুধ পৌঁছানোর ধারণা targeted drug delivery নামে পরিচিত।
\end{examplebox}

\subsubsection{ইলেকট্রনিক্সে ব্যবহার}

Computer chip, memory device, display, sensor, battery, solar cell এবং flexible electronics-এ nanoscale material গুরুত্বপূর্ণ। ক্ষুদ্র transistor ও efficient material electronic device-কে ছোট, দ্রুত ও energy-efficient করতে সাহায্য করে।

\subsubsection{কৃষি ও শিল্পে ব্যবহার}

কৃষিতে nano fertilizer, nano pesticide, soil sensor ও food packaging গবেষণায় ব্যবহৃত হয়। শিল্পে nano coating, stronger composite material, water filter, stain-resistant fabric এবং corrosion protection-এ ব্যবহার আছে।

\begin{center}
\begin{tabular}{|p{0.28\textwidth}|p{0.58\textwidth}|}
\hline
\textbf{ক্ষেত্র} & \textbf{ব্যবহার} \\
\hline
চিকিৎসা & drug delivery, diagnostic sensor, imaging \\
\hline
ইলেকট্রনিক্স & chip, memory, display, sensor \\
\hline
কৃষি & smart fertilizer, soil sensor, food packaging \\
\hline
শিল্প & nano coating, composite material, water purification \\
\hline
\end{tabular}
\end{center}

\subsubsection{ন্যানো টেকনোলজির সুবিধা}

\begin{itemize}
\item ক্ষুদ্র ও শক্তিশালী device তৈরি করা যায়।
\item material-এর strength, durability বা conductivity বাড়ানো যায়।
\item চিকিৎসায় targeted therapy ও দ্রুত diagnosis-এ সহায়তা করতে পারে।
\item পানি বিশুদ্ধকরণ ও পরিবেশ পর্যবেক্ষণে ব্যবহার করা যায়।
\item energy storage ও solar technology উন্নত করতে পারে।
\end{itemize}

\subsubsection{ন্যানো টেকনোলজির ঝুঁকি}

ন্যানো উপাদানের size খুব ছোট হওয়ায় শরীর, পানি, মাটি বা বায়ুতে এর আচরণ সাধারণ পদার্থের চেয়ে ভিন্ন হতে পারে। তাই toxicity, environmental impact, worker safety, regulation এবং ethical use নিয়ে গবেষণা ও সতর্কতা প্রয়োজন।

\begin{mistakebox}{সতর্কতা}
ন্যানো টেকনোলজি মানেই শুধু “ছোট জিনিস” নয়; nanoscale-এ পদার্থের বৈশিষ্ট্য বদলে যাওয়াই মূল বিষয়।
\end{mistakebox}

\begin{boardbox}{বোর্ড ফোকাস}
ন্যানো টেকনোলজির সংজ্ঞায় 1-100 nanometer scale উল্লেখ করলে উত্তর শক্তিশালী হয়। ব্যবহারক্ষেত্রে চিকিৎসা, ইলেকট্রনিক্স, কৃষি ও শিল্প থেকে উদাহরণ দাও।
\end{boardbox}

\begin{practicebox}
\textbf{MCQ}
\begin{enumerate}
\item ন্যানো টেকনোলজি সাধারণত কোন scale নিয়ে কাজ করে?\\
ক. 1-100 nanometer \quad খ. 1-100 kilometer \quad গ. 1-100 liter \quad ঘ. 1-100 byte only
\item Carbon nanotube কোন ধরনের উপাদানের উদাহরণ?\\
ক. Nanomaterial \quad খ. Spreadsheet \quad গ. Antivirus \quad ঘ. Router
\end{enumerate}

\textbf{সংক্ষিপ্ত প্রশ্ন}
\begin{enumerate}
\item ন্যানো টেকনোলজি কী?
\item চিকিৎসায় ন্যানো টেকনোলজির দুটি ব্যবহার লেখ।
\end{enumerate}
\end{practicebox}

\begin{sourcebox}
এই অংশে local ICT PDF resources, National Nanotechnology Initiative, Britannica nanotechnology overview, NIH nanomedicine articles এবং HSC ICT topic outline মিলিয়ে লেখা হয়েছে।
\end{sourcebox}
